
\section{Method}

\subsection{Overview}

SVD-Linear Attention replaces MHSA with a kernelized linear attention 
whose projections are SVD-factorized.

\textbf{Figure 4.1} shows the pipeline.

\begin{lstlisting}
    +--------------------------------------------+
    |  Input X $\in$ R^{B$\times$L$\times$d}                        |
    +----------------+---------------------------+
                     |
        +------------+------------+
                                   
    +---------------+      +---------------+
    | SVDLinear Q,K,V|      | decay\_gate    |
    | (factorized)  |      | (x $\to$ $\lambda$)       |
    +-------+-------+      +-------+-------+
            |                      |
                                    
    +---------------+      +---------------+
    |  (Q),  (K)    |      | Chunked scan  |
    | feature map   |      | S\_t, Z\_t      |
    +-------+-------+      +-------+-------+
            |                      |
            +----------+-----------+
                        
              +---------------+
              |  (Q) S /  (Q) Z|
              +-------+-------+
                       
              +---------------+
              | out\_proj (SVD)|
              +---------------+
\end{lstlisting}

\subsection{SVD-Factorized Projections}

Each projection is stored as $U V^\top$ where $U \in \mathbb{R}^{d_\text{out} \times r}$ 
and $V \in \mathbb{R}^{r \times d_\text{in}}$. The forward computes:

\[
y = (x V^\top) U^\top + b + \text{LoRA}(x)
\tag{4.1}
\]

where LoRA provides a trainable residual:

\[
\text{LoRA}(x) = (x A^\top) B^\top \cdot s
\tag{4.2}
\]

\textbf{Parameter count:} $r(d_\text{in} + d_\text{out})$ vs $d_\text{in} d_\text{out}$. 
For $768 \times 768$, $r = 32$: \textbf{~11$\times$ reduction}.

\textbf{Initialization.} $U$ and $V$ are initialized orthogonally.

\textbf{LoRA.} $A$ is Kaiming-uniform, $B = 0$.

\subsection{Kernelized Linear Attention}

\subsubsection{Feature Map}

\[
\phi(x) = \text{ELU}(x) + 1
\tag{4.3}
\]

which is strictly positive, ensuring the denominator $\phi(Q_t) Z_t > 0$.

\subsubsection{Bidirectional (Encoder) Mode}

For a non-causal encoder:

\[
S = \phi(K)^\top V, \qquad
Z = \sum\_j \phi(K\_j)
\tag{4.4}
\]

\[
\text{Out}\_t = \frac{\phi(Q\_t) S}{\phi(Q\_t) Z + \epsilon}
\tag{4.5}
\]

\textbf{Complexity:} $\mathcal{O}(N d_h^2)$.

\subsubsection{Causal (Decoder) Mode with Chunked Scan}

For a causal decoder with decay gating, we use a chunked scan (chunk size $C = 64$).

\textbf{Within chunk $c$:}

\[
\log \gamma\_{c, i} = \sum\_{j \leq i} \log \lambda\_{c, j}
\tag{4.6}
\]

\[
S\_{c, i}^{\text{intra}} = e^{\log \gamma\_{c, i}} \sum\_{j \leq i} e^{-\log \gamma\_{c, j}} \phi(K\_{c, j}) V\_{c, j}^\top
\tag{4.7}
\]

\textbf{Inter-chunk:}

\[
S\_{c, i} = S\_{c, i}^{\text{intra}} + e^{\log \gamma\_{c, i}} S\_{c-1, C}
\tag{4.8}
\]

\textbf{Output:}

\[
\text{Out}\_{c, i} = \frac{\phi(Q\_{c, i}) S\_{c, i}}{\phi(Q\_{c, i}) Z\_{c, i} + \epsilon}
\tag{4.9}
\]

\textbf{Complexity:} $\mathcal{O}(N \cdot C \cdot d_h)$ for the intra-chunk matrix, 
$\mathcal{O}(N d_h^2)$ for the state updates.

\subsection{Exponential Decay Gate}

\subsubsection{Decay Computation}

\[
\lambda\_t = \exp(-\text{softplus}(W\_g x\_t + b\_g))
\tag{4.10}
\]

For channel-wise decay:

\[
\lambda\_{t, h, k} = \exp(-\text{softplus}(W\_g x\_t)\_{h, k} + b\_{g, h, k})
\tag{4.11}
\]

\subsubsection{Initialization Strategy}

We want the initial decay times $\tau_{h} = -\log \lambda_h$ to be log-spaced:

\[
\tau\_h = \exp\!\left(\text{linspace}(\log 1, \log 1024, H)\right) \quad \text{for } h = 1, \dots, H
\tag{4.12}
\]

Since $\lambda = e^{-\tau}$ and $\lambda = e^{-\text{softplus}(b)}$, we set:

\[
b\_{g, h} = -\log \tau\_h = \text{softplus}^{-1}(\log \tau\_h)
\tag{4.13}
\]

In practice, we approximate $\text{softplus}(b) \approx \log(1 + e^b) \approx \log \tau_h$ for large $\tau_h$, so:

\[
b\_{g, h} \approx \log(\log \tau\_h)
\tag{4.14}
\]

The code uses \texttt{bias\_target = -torch.log(decay\_times)} which is the 
\textbf{exact} target for the \textbf{unbounded} log-decay, not the softplus-bounded one. The actual decay is:

\[
\lambda = e^{-\text{softplus}(W\_g x + b\_g)}
\]

so the effective initial decay is \textbf{not exactly} the target. We do not 
correct for this.

\subsubsection{Padding Handling}

For padding tokens, we force $\lambda_t = 1$:

\[
\lambda\_t = \begin{cases} 1 & \text{if } x\_t \text{ is padding} \\ \lambda\_t & \text{otherwise} \end{cases}
\tag{4.15}
\]

This ensures that padding does not artificially decay the memory state.

\subsection{Streaming Inference}

The \texttt{step\_recurrent} method processes one token at a time:

\[
S\_t = \lambda\_t S\_{t-1} + \phi(K\_t) V\_t^\top, \qquad
\text{Out}\_t = \frac{\phi(Q\_t) S\_t}{\phi(Q\_t) Z\_t + \epsilon}
\tag{4.16}
\]

\textbf{Complexity:} $\mathcal{O}(d_h^2)$ per step.

This is equivalent to the batch chunked scan when tokens are processed in 
order, but reduces memory from $\mathcal{O}(N d_h^2)$ to $\mathcal{O}(d_h^2)$.

\subsection{Full Algorithm}

\textbf{Algorithm 1:} SVD-Linear Attention batch forward pass
\begin{lstlisting}
Input:  X $\in$ R^{B$\times$N$\times$d}, attention\_mask
Output: Y $\in$ R^{B$\times$N$\times$d}

 1: Q $\leftarrow$ SVDLinear\_Q(X); K $\leftarrow$ SVDLinear\_K(X); V $\leftarrow$ SVDLinear\_V(X)
 2: if RoPE: (Q, K) $\leftarrow$ Rotary(Q, K)
 3: if GQA:  K, V $\leftarrow$ repeat\_kv(K, V)
 4:  \_Q $\leftarrow$ ELU(Q) + 1;  \_K $\leftarrow$ ELU(K) + 1
 5: apply padding mask to  \_K
 6: $\lambda$ $\leftarrow$ exp(-softplus(decay\_gate(X)))
 7: force $\lambda$=1 on padding
 8:
 9: # Chunked scan
10: for c in chunks of size C=64:
11:     compute intra-chunk scan
12:     add inter-chunk carry
13:     compute output for chunk c
14:     update running state
15:
16: Y $\leftarrow$ SVDLinear\_O(concat(outputs))
17: return Y
\end{lstlisting}

\textbf{Algorithm 2:} \texttt{step\_recurrent} (streaming)
\begin{lstlisting}
Input:  x\_t $\in$ R^{B$\times$d}, past\_state (S, Z)
Output: y\_t $\in$ R^{B$\times$d}, new\_state

1: q $\leftarrow$ SVDLinear\_Q(x\_t); k $\leftarrow$ SVDLinear\_K(x\_t); v $\leftarrow$ SVDLinear\_V(x\_t)
2:  \_q $\leftarrow$ ELU(q) + 1;  \_k $\leftarrow$ ELU(k) + 1
3: $\lambda$ $\leftarrow$ exp(-softplus(decay\_gate(x\_t)))
4: S $\leftarrow$ $\lambda$   S +  \_k^T v
5: Z $\leftarrow$ $\lambda$   Z +  \_k^T
6: num $\leftarrow$  \_q S
7: den $\leftarrow$  \_q Z + $\epsilon$
8: context $\leftarrow$ num / den
9: y\_t $\leftarrow$ SVDLinear\_O(context)
10: return y\_t, (S, Z)
\end{lstlisting}

\subsection{Complexity Analysis}

\begin{table}[H]
\centering
\small
\begin{tabular}{lll}
\toprule
\textbf{Mode} & \textbf{Time} & \textbf{Space} \\
\midrule
Softmax (baseline) & $\mathcal{O}(N^2 d_h)$ & $\mathcal{O}(N^2)$ \\
SVD-Linear (bidirectional) & $\mathcal{O}(N d_h^2)$ & $\mathcal{O}(d_h^2)$ \\
SVD-Linear (chunked causal) & $\mathcal{O}(N C d_h + N d_h^2)$ & $\mathcal{O}(C d_h^2)$ \\
SVD-Linear (step\_recurrent) & $\mathcal{O}(d_h^2)$ per step & $\mathcal{O}(d_h^2)$ \\
\bottomrule
\end{tabular}
\end{table}

For $N = 512$, $d_h = 64$, $C = 64$:

\begin{table}[H]
\centering
\small
\begin{tabular}{ll}
\toprule
\textbf{Mode} & \textbf{Time} \\
\midrule
Softmax & 512  $\times$ 64 $\approx$ 1.7 $\times$ 10 \\
SVD-Linear & 512 $\times$ 64 $\times$ 64 $\approx$ 2 $\times$ 10 \\
\textbf{Speedup} & \textbf{~8$\times$} \\
\bottomrule
\end{tabular}
\end{table}

\subsection{Parameter Analysis}

\begin{table}[H]
\centering
\small
\begin{tabular}{ll}
\toprule
\textbf{Component} & \textbf{Parameters} \\
\midrule
Q, K, V, O projections (SVD) & $4 \times r(d_\text{in} + d_\text{out})$ \\
LoRA adapters & $4 \times r_L(d_\text{in} + d_\text{out})$ \\
Decay gate & $d \times H + H$ \\
\textbf{Total} & see below \\
\bottomrule
\end{tabular}
\end{table}

For DistilBERT ($d = 768$, $H = 12$, $d_h = 64$, $r = 32$, $r_L = 16$, $N_\ell = 6$):

\begin{table}[H]
\centering
\small
\begin{tabular}{lll}
\toprule
\textbf{Component} & \textbf{Per layer} & \textbf{6 layers} \\
\midrule
Q/K/V/O SVD & 4 $\times$ 49,216 = 196,864 & 1,181,184 \\
LoRA & 4 $\times$ 24,576 = 98,304 & 589,824 \\
Decay gate & 9,228 & 55,368 \\
\textbf{Total} & \textbf{304,396} & \textbf{1,826,376} \\
\bottomrule
\end{tabular}
\end{table}

Baseline per layer: 2,359,296. \textbf{Reduction: ~50\%$\times$}.

\subsection{Theoretical Notes}

\subsubsection{Linear Attention Approximation}

\textbf{Claim.} The linear attention with $\phi(x) = \text{ELU}(x) + 1$ approximates 
softmax attention but does not equal it.

\textbf{Proof.} Linear attention is exact for $\phi(x) = e^x$, which is 
numerically unstable. For ELU+1, the approximation quality depends on the 
distribution of Q, K. In practice it gives similar top-$k$ attention patterns. $\square$

\subsubsection{Causality with Decay}

\textbf{Claim.} SVD-Linear Attention preserves causality.

\textbf{Proof.} The chunked scan is causal: chunk $c$ can only attend to chunks 
$1..c$ and intra-chunk positions $1..i$. The step-recurrent mode is causal 
by construction. $\square$

\subsubsection{Streaming Equivalence}

\textbf{Claim.} \texttt{step\_recurrent} produces the same output as batch chunked scan.

\textbf{Proof.} Both compute the same recurrence $S_t = \lambda_t S_{t-1} + \phi(K_t) V_t^\top$. 
The batch version uses a parallel scan; the streaming version is sequential. 
The results are mathematically identical. $\square$
\begin{lstlisting}

---
\end{lstlisting}